% Einheit 4 von 4 – Einheit 4 (bank/flaecheninhalt-und-volumen-im-raum/e4.jsonl, zusammenbau v0.1)
%% TODO Prüfungswort Einheit 4: nicht in der Marken-Zeile
%% TODO Zeitmarke Einheit 4: Marken-Zeile nicht lesbar
%% TODO Zweigzeile Teil 1: was hier gelernt wird (Satz in Schülersprache aus dem Einheitstitel) – Einheit 4
%% TODO Einheitentitel 4 nicht in der Mappe lesbar
\einheitenkopf[e4]{Einheit 4 von 4 · Einheit 4}
\setcounter{aufgabe}{16}

%% TODO Titel: Ich-kann-Satz fehlt in der Bank (Platzhalter: Kettenname) – Nr. 17
%% TODO Anweisung: ein Satz über der ersten Teilaufgabe fehlt in der Bank – Nr. 17
\begin{aufgabe}{Zusammengesetzt und Verhältnisse}
%% TODO \rechenplatz{n}: Schrittzahl des Grundfalls fehlt in der Bank (nur bei fester Schreibform, 2.3 b)
\begin{teile}
\teil Ein Haus: Quader mit aufgesetztem Satteldach. Wie fasst du sein Volumen? \\ \kreuz{direkt} \\ \kreuz{als Summe} \\ \kreuz{als Differenz}
\teil Ein Quader hat die Grundfläche $PQRU$ mit $P(1 | 1 | 0)$, $Q(7 | 1 | 0)$, $R(7 | 7 | 0)$, $U(1 | 7 | 0)$ und die Höhe 4. Die Ebene durch $Q$, $U$ und $T(1 | 1 | 8)$ zerlegt den Quader. Berechne das Volumen des Teilkörpers, der $P$ enthält, und erläutere den Ansatz. V = \leerfeld VE
\teil Ein Würfel mit der Kantenlänge 6 hat eine Ecke im Ursprung, seine Kanten liegen auf den positiven Achsen. Die Ebene durch die Würfelkante auf der $x$-Achse und den Mittelpunkt $M(6 | 3 | 6)$ einer oberen Kante teilt den Würfel. Volumen des kleineren Teilkörpers? \hfill (Abitur 2025 LK) \\ V = \leerfeld VE
\teil Ein Sockel hat die Form eines Stumpfs einer geraden quadratischen Pyramide: Grundfläche mit der Seite 6\,m, Deckfläche mit der Seite 3\,m, Höhe 6\,m. Sein Volumen wird mit dem Term $\tfrac{1}{3} \cdot 36 \cdot 12 - \tfrac{1}{3} \cdot 9 \cdot 6$ berechnet. Begründe, dass der Term das Volumen liefert. \hfill (Abitur 2022 GK)
\teil Ein Quader und eine Pyramide haben dieselbe Grundfläche, die Spitze der Pyramide liegt in der Deckfläche des Quaders. Um wie viel Prozent ist das Volumen des Quaders größer als das der Pyramide? Rechne ohne konkrete Volumina. \hfill (Abitur 2021 GK) \\ \leerfeld \%
\teil Das Dreieck mit den Ecken $(0 | 0)$, $(8 | 0)$ und $(8 | 6)$ rotiert um die $x$-Achse. Oberflächeninhalt des entstehenden Körpers? \hfill (Abitur 2018 GK) \\ O = \leerfeld FE
\teil Eine zylindrische Vorratsdose mit dem Durchmesser 20\,cm fasst 10 Liter. Passt sie in ein Regalfach mit 30\,cm lichter Höhe? \hfill (Abitur 2022 GK)
\teil Die Pyramide mit den Ecken $O$, $P(6 | 0 | 0)$, $Q(0 | 7 | 0)$ und $R(0 | 0 | 5)$ wird von einer Ebene parallel zur $y$-$z$-Ebene in zwei volumengleiche Teile zerlegt. An welcher Stelle $t$ schneidet die Ebene die $x$-Achse? \hfill (Abitur 2022 LK) \\ t = \leerfeld
\end{teile}
\end{aufgabe}

%% TODO Titel: Ich-kann-Satz fehlt in der Bank (Platzhalter: Kettenname) – Nr. 18
%% TODO Anweisung: ein Satz über der ersten Teilaufgabe fehlt in der Bank – Nr. 18
\begin{aufgabe}{Zusammengesetzt und Verhältnisse – Fehler finden, Begründen, Darstellungswechsel, Anwendung}
\begin{teile}
\teil Ein Quader hat eine quadratische Grundfläche mit der Seite 6 und die Höhe 3. Die Ebene durch zwei Nachbarecken $Q$ und $U$ der Grundfläche und den Punkt $T$ senkrecht über ihrer gemeinsamen Nachbarecke $P$ in der Höhe 9 zerlegt den Quader. Paul rechnet für den Teilkörper mit $P$: \rechnung{V &= \tfrac{1}{3} \cdot \tfrac{1}{2} \cdot 6 \cdot 6 \cdot 9 = 54} Finde den Fehler und rechne richtig.
\teil Begründe: Schneidet eine Ebene parallel zur Grundfläche eine Teilpyramide ab, so ist ihr Volumenanteil der Kubus des Streckfaktors.
\teil Das Volumen eines Restkörpers ist $V(t) = 40 - 5t$ für $0 \le t \le 6$. Zeichne den Graphen von $V$ in das Koordinatensystem.

\begin{ksys}[xmin=0,xmax=7,ymin=0,ymax=45,xstep=1,ystep=5,xlabel=t,ylabel=V] \end{ksys}
\teil Ein Betonsockel ist der Stumpf einer geraden quadratischen Pyramide: unten Seite 1{,}2\,m, oben Seite 0{,}8\,m, Höhe 6\,m; die verlängerten Seitenkanten treffen sich 18\,m über dem Boden. Beton wiegt 2{,}4\,t je m³. Masse des Sockels? \hfill (Abitur 2018 LK) \\ \leerfeld[t]
\end{teile}
\end{aufgabe}
